\begin{proof}
The global nonemptiness hypothesis in
\Cref{prop:p033-invariant-toric-real-pl-envelope} restricts to a member of
\(\mathcal C_{\varphi,B}\).  Applying the local convex-envelope argument in
that proposition with the convex domain \(B\) shows that \(p_B\) is finite,
continuous, and convex.

If \(h\in\mathcal L^S_{f,B}\), then \cite[\S7.14]{GJR26} makes \(h\) locally
psh-approachable.  Since \(V_B\) is good, separated, strictly analytic, and
boundaryless, the same result makes \(h\) a finite continuous \(S\)-invariant
psh function.  The toric correspondence
\cite[Theorem~5.3.11]{BGJKK25} therefore writes
\(h=a\circ\operatorname{trop}\) for a finite convex function \(a\) on
\(B\).  Since \(h\leq f\), we have \(a\leq\varphi|_B\), whence
\[
                         h\leq p_B\circ\operatorname{trop}.
\tag{1}\label{eq:p033-lfs-upper}
\]

We construct the reverse approximation directly.  Choose an increasing compact
exhaustion \((K_j)_j\) of \(B\) with \(K_j\Subset B\); for example, after
discarding empty initial terms, take
\[
 K_j:=\{x\in B:\operatorname{dist}(x,\mathbf R^n\setminus B)\geq j^{-1}\}.
\]
Set \(\epsilon_j=\eta_j=2^{-j}\).  At every point of \(K_j\), the finite convex
function \(p_B\) has a supporting affine function.  The open sets on which
one such support is within \(\epsilon_j\) of \(p_B\) cover \(K_j\).  A finite
subcover gives affine functions \(\lambda_{ij}\) such that
\[
 p_B-\epsilon_j\leq p_j:=\max_i\lambda_{ij}\leq p_B
 \qquad\hbox{on }K_j.
\tag{2}\label{eq:p033-lfs-finite-support}
\]

Write \(M=N^\vee\).  We next round every \(\lambda_{ij}\) from below on the
fixed compact set \(\overline B\).  Namely, for an affine
\(\lambda(x)=\langle a,x\rangle+b\) and \(\eta>0\), there are an integer
\(r\geq1\), an \(m\in M\), and a \(c\in-\log|K^\times|\) such that
\[
 0\leq\lambda(x)-r^{-1}\bigl(\langle m,x\rangle+c\bigr)\leq\eta
 \qquad\hbox{on }\overline B.
\tag{3}\label{eq:p033-lfs-lattice-minorant}
\]
Indeed, choose a nonzero \(\gamma\in-\log|K^\times|\), round \(ra\) to
\(m\), and choose \(c\in\gamma\mathbf Z\) just below the scalar that makes
the displayed affine difference nonnegative.  As \(r\) grows, the rounding
error and the mesh \(\lvert\gamma\rvert/r\) are both arbitrarily small.  One
common sufficiently large \(r_j\) works for the finite family in
\eqref{eq:p033-lfs-finite-support}.

Let \(L_{ij}=r_j^{-1}(\langle m_{ij},\cdot\rangle+c_{ij})\) be the resulting
minorants and put \(Q_j:=\max_iL_{ij}\).  Then \(Q_j\leq p_B\) on \(B\), and
\eqref{eq:p033-lfs-finite-support}--\eqref{eq:p033-lfs-lattice-minorant} give
\[
 p_B-\epsilon_j-\eta_j\leq Q_j\leq p_B
 \qquad\hbox{on }K_j.
\tag{4}\label{eq:p033-lfs-approximation}
\]
Choose \(a_{ij}\in K^\times\) with
\(-\log|a_{ij}|=c_{ij}\).  The torus convention gives
\[
 Q_j\circ\operatorname{trop}
 =\frac1{r_j}\max_i\{-\log|a_{ij}\chi^{m_{ij}}|\}.
\tag{5}\label{eq:p033-lfs-normal-form}
\]
Thus \(Q_j\circ\operatorname{trop}\) is locally Fubini--Study by
\cite[\S7.9]{GJR26}.  It is finite and \(S\)-invariant, and its upper bound
in \eqref{eq:p033-lfs-approximation} gives
\[
 Q_j\circ\operatorname{trop}
 \leq p_B\circ\operatorname{trop}
 \leq f|_{V_B}.
\]
Hence it belongs to \(\mathcal L^S_{f,B}\).  Since every compact subset of
\(B\) lies in \(K_j\) for all sufficiently large \(j\),
\eqref{eq:p033-lfs-approximation} proves locally uniform convergence from
below to \(p_B\circ\operatorname{trop}\).

The approximants can be made monotone without leaving the locally
Fubini--Study class.  Put
\[
                       R_j:=\operatorname{lcm}(r_1,\ldots,r_j),
 \qquad \widetilde Q_j:=\max_{1\leq\ell\leq j}Q_\ell .
\]
Writing each $Q_\ell$ in the normal form
\eqref{eq:p033-lfs-normal-form} and multiplying its finitely many affine
functions by the integer $R_j/r_\ell$ gives
\[
 \widetilde Q_j\circ\operatorname{trop}
 =\frac1{R_j}\max_{\substack{1\leq\ell\leq j\\ i}}
 \left\{-\log\left|a_{i\ell}^{R_j/r_\ell}
       \chi^{(R_j/r_\ell)m_{i\ell}}\right|\right\}.
\tag{6}\label{eq:p033-lfs-monotone-normal-form}
\]
Thus every \(\widetilde Q_j\circ\operatorname{trop}\) again belongs to
\(\mathcal L^S_{f,B}\).  The sequence is pointwise nondecreasing and,
because \(Q_j\leq\widetilde Q_j\leq p_B\), the estimate
\eqref{eq:p033-lfs-approximation} shows that it converges locally uniformly
to \(p_B\circ\operatorname{trop}\).  This supplies the reverse inequality
to \eqref{eq:p033-lfs-upper}, and proves the claim.
\end{proof}
